% sec-02: outline item B.  Figures 4 (diagrams), 5 (d2), 6 (plantuml).
\section{The Entity and the Asset}

\subsection{What the Entity Is}

ChemicalQDevice is a California corporation. Its sole officer, sole director and
sole owner is Kevin Kawchak. It has no employees at the date of writing, no
negotiated indirect-cost rate agreement, and no facilities of its own. It is
therefore more than 50 percent directly owned and controlled by one United
States individual, which is the eligibility test the SBIR programme applies
\cite{ecfr2026sbirownership}, and \S4 is written to keep it so through the
thirty-three months this plan covers.

\begin{appfloat}[!tb]
\begin{appfig}[diagrams-python-type / clustered zones]
% Four zones.  Tile pitch 27mm horizontal, 22mm vertical.  Every label is set
% 5.4mm beneath its tile by \dgnode, and every cluster fit names both the tile
% node and its label node, so no dashed border cuts a label.
\dgnodew{dgtiled}{0}{0}{\glyphdoc}{Phase 1 protocol}{o1}
\dgnode{dgtile}{2.7}{0}{\glyphdoc}{Phase 2 protocol}{o2}
\dgnodew{dgtiled}{5.4}{0}{\glyphflask}{IND package, drafted}{o3}
\dgnodew{dgtilem}{8.1}{0}{\glyphcpu}{QSP and VVUQ suite}{o4}
\dgnode{dgtile}{0}{-2.2}{\glyphdb}{Repository, MIT}{o5}
\dgnode{dgtile}{2.7}{-2.2}{\glyphbank}{Ten application sets}{o6}
\dgnode{dgtile}{5.4}{-2.2}{\glyphdoc}{Four legislative drafts}{o7}
\begin{scope}[on background layer]
\node[dgcluster,fit=(o1)(o1l)(o2)(o2l)(o3)(o3l)(o4)(o4l)(o5)(o5l)(o6)(o6l)(o7)(o7l),
      inner sep=7pt] (zo) {};
\end{scope}
\node[dgctitle,anchor=south west] at ([yshift=1.2mm]zo.north west)
  {Owned outright, no encumbrance};
\node[d2cellk,minimum width=13mm,text width=10mm,anchor=north east]
  at ([xshift=-1.5mm,yshift=-1.5mm]zo.north east) {13 rows};
\dgnodeg{dgtileg}{0}{-4.9}{\glyphai}{Base model weights}{l1}
\dgnodeg{dgtileg}{2.7}{-4.9}{\glyphrobot}{Robotic platform, site}{l2}
\dgnodeg{dgtileg}{5.4}{-4.9}{\glyphpill}{Daraxonrasib, developer}{l3}
\begin{scope}[on background layer]
\node[dgcluster2,fit=(l1)(l1l)(l2)(l2l)(l3)(l3l),inner sep=7pt] (zl) {};
\end{scope}
\node[dgctitle2,anchor=south west] at ([yshift=1.2mm]zl.north west)
  {Licensed, none exclusive};
\node[d2cellg,minimum width=13mm,text width=10mm,anchor=north east]
  at ([xshift=-1.5mm,yshift=-1.5mm]zl.north east) {3 rows};
\node[dgcluster2,minimum width=29mm,minimum height=19mm,anchor=north west] (zc) at (8.6,-4.05) {};
\node[dgctitle2,anchor=south west] at ([yshift=1.2mm]zc.north west) {Contracted};
\node[font=\tiny\sffamily\itshape,text=protogray,text width=24mm,align=center] at (zc.center)
  {no rows exist as of August 2026};
\dgnodeg{dgtilek}{0}{-7.6}{\glyphhand}{Site agreement}{a1}
\dgnodeg{dgtilek}{2.7}{-7.6}{\glyphshield}{IRB approval}{a2}
\dgnodeg{dgtilek}{5.4}{-7.6}{\glyphpill}{Drug supply}{a3}
\dgnodeg{dgtilek}{8.1}{-7.6}{\glyphlink}{Letter of authorization}{a4}
\begin{scope}[on background layer]
\node[dgcluster2,fit=(a1)(a1l)(a2)(a2l)(a3)(a3l)(a4)(a4l),inner sep=7pt] (za) {};
\end{scope}
\node[dgctitle2,anchor=south west] at ([yshift=1.2mm]za.north west)
  {Absent, none obtainable by the company alone};
\node[d2cellg,minimum width=13mm,text width=10mm,anchor=north east]
  at ([xshift=-1.5mm,yshift=-1.5mm]za.north east) {7 rows};
\draw[dgedgeb] (zo.south west) to[bend right=14]
  node[font=\tiny\sffamily,fill=protowhite,inner sep=1.5pt,pos=0.5] {supplies} (za.north west);
\draw[dgedged] (zl.south east) -- node[font=\tiny\sffamily,fill=protowhite,inner sep=1.5pt,pos=0.5]
  {depends on} (za.north east);
\node[font=\tiny\itshape,text=protogray,anchor=north west,text width=132mm]
  at (-1.3,-8.95) {The contracted zone is drawn at the size a three-tile zone
  would occupy, so its emptiness reads as absence at scale rather than as a
  small box. The pill glyph appears twice: the same object is licensed to
  somebody else and absent from this company.};
\end{appfig}
\vspace{-0.65cm}
\figcaption{Four zones and thirteen tiles. The owned zone holds everything the\\
company made; the contracted zone is drawn at full size and holds\\
nothing, because an empty zone is the finding this figure reports.}
\end{appfloat}

\subsection{Owned}

\begin{apptable}
\begin{tabularx}{\textwidth}{>{\raggedright\arraybackslash}X >{\raggedright\arraybackslash}p{1.9cm} >{\raggedright\arraybackslash}p{1.7cm} >{\raggedright\arraybackslash}p{4.4cm}}
\headrow \thead{Asset} & \thead{Instrument} & \thead{First deposit} & \thead{Identifier}\\
\midrule
Phase 1 protocol & Zenodo & 2026-06-21 & \dlink{10.5281/zenodo.20780121} \\
Phase 2 protocol & Zenodo & 2026-06-24 & \dlink{10.5281/zenodo.20807027} \\
IND package, drafted, not filed & Zenodo & 2026-07-01 & \dlink{10.5281/zenodo.21097442} \\
Phase 1 LLM document guidance & Zenodo & 2026-06-29 & \dlink{10.5281/zenodo.21018646} \\
PI LLM adoption guide & Zenodo & 2026-06-26 & \dlink{10.5281/zenodo.20843290} \\
QSP simulation stack and VVUQ suite & Zenodo & 2025-08 & \dlink{10.5281/zenodo.17001137} \\
Daraxonrasib identification meta-analysis & Zenodo & 2025-06 & \dlink{10.5281/zenodo.15735068} \\
Funding application v1.0 & Zenodo & 2026-07-07 & \dlink{10.5281/zenodo.21232965} \\
Funding application v2.0 & Zenodo & 2026-07-12 & \dlink{10.5281/zenodo.21317266} \\
Patient robot advocacy paper & Zenodo & 2026-07-31 & \dlink{10.5281/zenodo.21720120} \\
H. R. 9510 bill v5.0 and companions & Zenodo & 2026-06-12 & \dlink{10.5281/zenodo.20619762} \\
Repository, MIT licensed & GitHub & 2026-08 & \appfile{physical-ai-oncology-trials} v4.5.0 \\
Ten funding application file sets & Repository & 2026-08-04 & \appfile{funding/pdac-funding-applications} \\
\end{tabularx}
\end{apptable}
\tabcap{Thirteen owned assets, every one deposited and addressable. No row
carries an encumbrance, no row is subject to a licence granted to a third party,
and no row was produced under a grant that reserves rights in it.}

What the thirteen rows cost to produce is on the record and is unusual enough to
state. The empirical triplicate, the ten-arm quantitative systems pharmacology
virtual trial, and the digital-twin platform were each produced for an estimated
\$36,330 and in about one month, against industry benchmarks of above \$120,000,
above \$2,000,000, and \$28,000 to \$700,000 respectively \cite{qspmetpancre}.
Those figures are estimates from the source set and are reported as such, not as
audited accounts; Table 16 carries them with their benchmarks.

The consequence for a funder is that none of the thirteen appears in the budget.
The protocol is written, the simulation stack is built, the verification suite
is frozen, and the IND package is drafted \cite{foundingdocs,onpremwhippl,%
phase2proto,phase1ind}. A request for \$1,606,000 is a request to run a trial,
not to write one.

\begin{appfloat}[!tb]
\begin{appfig}[d2-type / sql tables]
% Four sql_table records in two columns, plus a milestone record set apart to
% the right.  Every field line is name : type; row counts sit in a separate
% header strip because a count is metadata about a table, not a column of it.
\node[d2sql,text width=34mm,anchor=north west] (rown) at (0,0)
  {\textbf{owned}\nodepart{two}asset\_id : int, pk\\title : text\\instrument :
   enum\\first\_deposit : date\\encumbrance : enum, none};
\node[d2sql,text width=34mm,anchor=north west] (rlic) at (0,-4.30)
  {\textbf{licensed}\nodepart{two}item : text\\licensor : text\\terms : text,
   public\\exclusive : bool, false\\survives\_stop : bool};
\node[d2sql,text width=38mm,anchor=north west,
      rectangle split part fill={protogray,pagrayl}] (rabs) at (5.30,0)
  {\textbf{\textcolor{protowhite}{absent}}\nodepart{two}item : text\\%
   blocked\_milestone : fk\\signatory : text\\obtainable\_alone : bool, false};
\node[d2sql,text width=34mm,anchor=north west,
      rectangle split part fill={protogray,pagrayl}] (rcon) at (5.30,-4.30)
  {\textbf{\textcolor{protowhite}{contracted}}\nodepart{two}row\_count : int, zero};
\node[d2sql,text width=34mm,anchor=north west,
      rectangle split part fill={pablue1,pablue2}] (rmil) at (10.60,-1.90)
  {\textbf{\textcolor{protowhite}{milestone}}\nodepart{two}milestone\_id : text,
   pk\\months : text\\cost\_usd : int\\artifact : text};
\foreach \n/\c/\sty in {rown/13/d2cellk, rlic/3/d2cellg, rabs/7/d2cellg, rcon/0/d2cellg}{
  \node[\sty,minimum width=12mm,text width=9mm,anchor=north east]
    at ([xshift=-1.5mm,yshift=-1.2mm]\n.north east) {\c\ rows};}
\draw[d2edgeb] (rabs.east) -- node[font=\tiny\sffamily,fill=protowhite,inner sep=1.5pt,pos=0.5]
  {foreign key} ([yshift=6mm]rmil.west);
\draw[d2edge] (rown.south) to[bend right=16]
  node[font=\tiny\sffamily,fill=protowhite,inner sep=1.5pt,pos=0.5] {supplies evidence for}
  ([yshift=-14mm]rmil.west);
\draw[d2edged] (rcon.north) -- node[font=\tiny\sffamily,fill=protowhite,inner sep=1.5pt,pos=0.5]
  {each row would close one} (rabs.south);
\node[umlnote,text width=40mm,anchor=north west] at (5.30,-5.55)
  {An empty table is a finding. It is drawn with its header and no body row,
  because a missing table would instead be an omission.};
\node[font=\tiny\itshape,text=protogray,anchor=north west,text width=132mm]
  at (0,-7.55) {No row of the owned table is encumbered and no row of the
  licensed table is exclusive, so nothing in either can be pledged. The
  capitalization plan in section 4 therefore prices equity in the company and
  never a right in an asset.};
\end{appfig}
\vspace{-0.65cm}
\figcaption{The asset register as four typed records. Thirteen rows are owned\\
outright, three are licensed on public terms, seven are absent,\\
and the contracted table has no rows at all as of August 2026.}
\end{appfloat}

\subsection{Licensed, and Contracted}

Three things are licensed and none exclusively. The base model weights are used
under their vendor's public terms. The robotic platform is the site's asset and
not the company's. Daraxonrasib is the developer's investigational agent, and no
licence of any kind exists. None of the three survives a stop, because none is
owned.

The contracted table has no rows. No clinical trial agreement, no drug supply
agreement, no letter of authorization, no confidentiality agreement, no
subaward, and no trial-specific insurance binder is in force.

\subsection{Absent}

\begin{apptable}
\begin{tabularx}{\textwidth}{>{\raggedright\arraybackslash}X >{\raggedright\arraybackslash}p{1.6cm} >{\raggedright\arraybackslash}p{3.4cm} >{\raggedright\arraybackslash}p{3.4cm}}
\headrow \thead{Instrument absent} & \thead{Blocks} & \thead{Only signatory} & \thead{Obtainable by the company alone}\\
\midrule
Site agreement, executed & M1 & UC San Diego & No \\
IRB approval & M2 & Site institutional review board & No \\
Drug supply agreement & M6 & Revolution Medicines & No \\
Letter of authorization, IND cross-reference & M5 & Revolution Medicines & No \\
Theatre and robotic platform time & M7 & UC San Diego & No \\
Trial-specific liability insurance & M6 & Carrier & No, needs an executed CTA \\
A second full-time employee & M3 & Funded by the Phase I award & No, needs the award \\
\end{tabularx}
\end{apptable}
\tabcap{Seven absent instruments, the milestone each blocks, and the single
party that can supply it. The last column reads no on every row, which is the
whole content of the table and the reason \S8 exists.}

\begin{appfloat}[!tb]
\begin{appfig}[plantuml-type / use case with boundaries]
% Three system boundaries separated by 9mm corridors that carry no node.  Only
% association lines cross a corridor.  Actors sit at the two margins and never
% between boundaries.
\umlactor{ceo}{-0.75}{1.15}{Sponsor CEO}
\umlactor{pi}{-0.75}{-2.85}{Site PI}
\umlactor{irb}{-0.75}{-4.75}{Site IRB}
\umlactor{fda}{14.35}{1.15}{FDA}
\umlactor{dev}{14.35}{-6.95}{Drug developer}
\foreach \i/\x/\y/\lab in {1/2.30/1.15/{Author and amend the protocol},
  2/5.55/1.15/{Hold and maintain the IND}, 3/8.80/1.15/{Build and verify the interlock rig},
  4/2.30/-0.20/{Freeze and hash the VVUQ suite}, 5/5.55/-0.20/{File expedited safety reports},
  6/8.80/-0.20/{Deposit the public archive}}{
  \node[umlusecase,text width=25mm] (s\i) at (\x,\y) {\lab};}
\begin{scope}[on background layer]
\node[umlpkg,fit=(s1)(s2)(s3)(s4)(s5)(s6),inner sep=6pt] (bs) {};
\end{scope}
\node[umlpkgtab,anchor=south west] at ([yshift=0.6mm]bs.north west)
  {ChemicalQDevice, sponsor scope};
\foreach \i/\x/\y/\lab in {1/2.30/-2.85/{Consent a participant},
  2/5.55/-2.85/{Assign a dose level}, 3/8.80/-2.85/{Perform the operation},
  4/2.30/-4.20/{Adjudicate an endpoint}, 5/5.55/-4.20/{Grade a fistula, ISGPS},
  6/8.80/-4.20/{Convene the DSMB}}{
  \node[umlusecase,text width=25mm] (u\i) at (\x,\y) {\lab};}
\begin{scope}[on background layer]
\node[umlpkg,fit=(u1)(u2)(u3)(u4)(u5)(u6),inner sep=6pt] (bu) {};
\end{scope}
\node[umlpkgtab,anchor=south west] at ([yshift=0.6mm]bu.north west)
  {UC San Diego Moores, site scope};
\foreach \i/\x/\lab in {1/2.30/{Supply investigational drug},
  2/6.15/{Schedule theatre and robotic time}, 3/10.00/{Cross reference the developer IND}}{
  \node[umlusegray,text width=27mm] (b\i) at (\x,-6.95) {\lab};}
\begin{scope}[on background layer]
\node[umlpkg,fill=pagrayl,fit=(b1)(b2)(b3),inner sep=6pt] (bb) {};
\end{scope}
\node[umlpkgtab,fill=pagraym,anchor=south west] at ([yshift=0.6mm]bb.north west)
  {Blocked, no instrument exists};
\foreach \t in {s1,s2,s3}{\draw[umlassoc] (ceo) -- (\t);}
\foreach \t in {u1,u2,u3}{\draw[umlassoc] (pi) -- (\t);}
\draw[umlassoc] (irb) -- (u4);
\draw[umlassoc] (irb) -- (u6);
\draw[umldash] (fda) -- (s2);
\draw[umldash] (fda) -- (s5);
\draw[umldash] (dev) -- (b1);
\draw[umldash] (dev) -- (b3);
\node[umlnote,text width=34mm,anchor=north west] at (11.20,-2.85)
  {Each blocked case needs exactly one signature: a supply agreement, an
   executed CTA, and a letter of authorization.};
\node[font=\tiny\itshape,text=protogray,anchor=north west,text width=132mm]
  at (-1.6,-8.35) {The two corridors between boundaries carry no node at all.
  The only ink in them is the association lines that cross, which is what makes
  the boundaries read as boundaries.};
\end{appfig}
\vspace{-0.65cm}
\figcaption{Two system boundaries, twelve use cases, and the corridor between\\
them holding three that neither party may execute today. Each of\\
the three is blocked by exactly one missing signed instrument.}
\end{appfloat}

Three use cases sit in the corridor, and each is blocked by exactly one
signature. No amount of company work removes any of the three, which is the
honest reading of the register and the reason the twelve milestones in \S5 open
with a site feasibility memorandum rather than with an experiment.

The register is dated. On 13 January 2026, when the first proposal was released,
the owned table held two rows and the absent table held the same seven
\cite{paiproposal}. In the seven months since, the owned table has grown by
eleven rows and the absent table has not lost one. That asymmetry is the single
most useful thing in this section: it says precisely which kind of work the
author can do alone and which kind he cannot.
